% Theorem J: the finite-m closed loop as a perturbed gradient flow (PARADIGM.md, section 6d).
\section{Theorem J: the finite-$m$ closed loop as a perturbed gradient flow}\label{sec:thmJ}

This section proves the first step of Section~\ref{sec:program}, branch 1: at finite $m$, without dither, in the
scaling $\sigma = \sigma_0/k_e$, the real closed loop has, near every hyperbolic critical point of the $m$-particle
regularised occupation loss, a unique periodic orbit with the same stability type, at distance $O(1/k_e)$, whose
first-order location is the one Proposition H.7 and Remark H.7a predict. It turns Remark H.7a's ``formal''
second-order averaging into a theorem for this system.

\paragraph{Setting.} The closed loop of Theorem H's design with $\tau = 0$, at finite $m$, in the scaling
$\sigma = \sigma_0/k_e$:
\begin{equation}\label{eq:thmJ:CL}
  \begin{gathered}
    \dot x = f(x) + \dot x_m - k_e e - \fhat_\Theta(x),\qquad e := x - x_m(t),\qquad
    \fhat_\Theta := \frac1m\sum_j\tilde\Phi(\cdot\,;\vartheta_j),\\
    \dot\vartheta_i = G\nabla_\vartheta\bigl[e^\T\tilde\Phi(x;\vartheta_i)\bigr]
    - \frac{\sigma_0}{k_e}(\vartheta_i - \vartheta^0),\qquad i = 1,\dots,m,
  \end{gathered}
  \tag{CL}
\end{equation}
with $\Theta := (\vartheta_1,\dots,\vartheta_m)\in\R^{mp}$, $G_m := I_m\otimes G$, saturated features, $x_m\in C^2$
$T_0$-periodic with occupation measure $\rho$, and $f\in C^2$ on a neighbourhood of the tube
$\mathcal T_R := \{x : \operatorname{dist}(x,x_m([0,T_0])) \le R\}$ for some $R>0$. Put
\begin{gather*}
  r_\Theta(x) := f(x) - \fhat_\Theta(x),\qquad
  J_m(\Theta) := m\bar L_m(\Theta) + \frac{\sigma_0}{2}\sum_i\norm{\vartheta_i - \vartheta^0}^2_{\Ginv},\\
  \bar L_m(\Theta) := \frac12\int\norm{r_\Theta}^2\dd\rho,\\
  H_i(\Theta,t) := G\nabla_\vartheta\bigl[r_\Theta(x_m(t))^\T\tilde\Phi(x_m(t);\vartheta)\bigr]\Big|_{\vartheta=\vartheta_i}
  - \sigma_0(\vartheta_i - \vartheta^0),\\
  \langle H\rangle(\Theta) = -G_m\nabla J_m(\Theta).
\end{gather*}
The averaged system of \textup{(CL)} is
$\dot\Theta = \frac1{k_e}\langle H\rangle(\Theta) = -\frac1{k_e}G_m\nabla J_m(\Theta)$: the gradient flow of $J_m$ in
the $G^{-1}$ metric, which is $\Aavg$ at $\tau = 0$ for the $m$-atom measure (Remark H.0); in slow time
$s := t/k_e$ it does not depend on $k_e$. Its equilibria are the critical points of $J_m$. A critical point $\Theta^*$
is \emph{hyperbolic} if $\nabla^2J_m(\Theta^*)$ is non-singular; its index is the number of negative eigenvalues; it
is a \emph{non-degenerate minimum} if $\nabla^2J_m(\Theta^*)\succ0$. Write $\nu_1\le\dots\le\nu_{mp}$ for the
eigenvalues of $G_m^{1/2}\nabla^2J_m(\Theta^*)G_m^{1/2}$ (those of $G_m\nabla^2J_m(\Theta^*)$), and
$\lambda_* := \nu_1$ when $\Theta^*$ is a minimum.

Let $u(\Theta,t)$ be the zero-mean antiderivative in $t$ of $H(\Theta,t) - \langle H\rangle(\Theta)$
($T_0$-periodic), and define, as in Proposition H.7 and Remark H.7a, at the atomic measure
$\mu^* := \frac1m\sum_j\delta_{\vartheta^*_j}$:
\begin{gather*}
  J := D_\Theta\bigl[\tfrac1{k_e}\langle H\rangle\bigr](\Theta^*) = -\frac1{k_e}G_m\nabla^2J_m(\Theta^*),\\
  b_{\mathrm{frozen},i} := \frac1{k_e^2}G\nabla_\vartheta\bigl(U_{\mathrm{even}} + U_{\mathrm{odd}}\bigr)[\mu^*](\vartheta)
  \Big|_{\vartheta=\vartheta^*_i},\qquad
  b_{\mathrm{osc}} := \frac1{k_e^2}\bigl\langle D_\Theta H(\Theta^*,t)\cdot u(\Theta^*,t)\bigr\rangle_t.
\end{gather*}

\begin{thm}{Theorem}{J}
Let $\Theta^*$ be a hyperbolic critical point of $J_m$. There are $k_0,\rho_0,C>0$, depending on \textup{(}$f$, $x_m$,
$G$, $\sigma_0$, $R$, the feature bounds, $\Theta^*$\textup{)} but not on $k_e$, such that for every $k_e\ge k_0$:
\begin{enumerate}[label=\textup{(\alph*)}]
  \item \emph{Existence, uniqueness, location.} \textup{(CL)} has a $T_0$-periodic solution $(x_p,\Theta_p)$, unique
    among the $T_0$-periodic solutions with $\sup_t\norm{\Theta_p(t) - \Theta^*}_{\Ginv} \le \rho_0$ and
    $\sup_t\norm{x_p(t) - x_m(t)} \le R$. It satisfies
    \begin{gather*}
      \sup_t\norm{x_p(t) - x_m(t)} \le C/k_e,\qquad \sup_t\norm{\Theta_p(t) - \Theta^*}_{\Ginv} \le C/k_e,\\
      \Theta_p(t) - \bar\Theta_p = \frac1{k_e}u(\Theta^*,t) + O(k_e^{-2}),
    \end{gather*}
    with $\bar\Theta_p$ the period mean of $\Theta_p$.
  \item \emph{Stability type.} The monodromy matrix of $(x_p,\Theta_p)$ has $n$ eigenvalues of modulus at most
    $e^{-(k_e - L)T_0}$, with $L$ independent of $k_e$ \textup{(}the error channel\textup{)}, and $mp$ eigenvalues
    \[
      \lambda_j = \exp\bigl(-(T_0/k_e)\nu_j + O(k_e^{-2})\bigr),\qquad j = 1,\dots,mp.
    \]
    Hence the orbit is hyperbolic with the index of $\Theta^*$: exponentially stable if $\nabla^2J_m(\Theta^*)\succ0$,
    with parameter Floquet exponents $(1/k_e)\nu_j + O(k_e^{-2})$, and unstable otherwise.
  \item \emph{Basin \textup{(}non-degenerate minimum\textup{)}.} If $\nabla^2J_m(\Theta^*)\succ0$, every solution of
    \textup{(CL)} with $\norm{\Theta(0) - \Theta^*}_{\Ginv} \le \rho_0$ and $\norm{x(0) - x_m(0)} \le R$ converges to
    the periodic orbit, and
    \begin{gather*}
      \norm{\Theta(t) - \Theta_p(t)}_{\Ginv} + \norm{x(t) - x_p(t)}
      \le Ce^{-\mu t}\bigl(\norm{\Theta(0) - \Theta_p(0)}_{\Ginv} + \norm{x(0) - x_p(0)}\bigr),\\
      \mu := \frac{\lambda_*}{k_e}(1 - \kappa),
    \end{gather*}
    with $\kappa = \kappa(\rho_0,1/k_e)\to0$ as $(\rho_0,1/k_e)\to0$. The basin radius $\rho_0$ does not depend on
    $k_e$; the rate is $\lambda_*/k_e$ to leading order, that is, $\lambda_*$ in slow time.
  \item \emph{First-order location.} The period mean satisfies
    \begin{align*}
      \bar\Theta_p &= \Theta^* - J^{-1}(b_{\mathrm{frozen}} + b_{\mathrm{osc}}) + O(k_e^{-2})\\
      &= \Theta^* + \frac1{k_e}\bigl(G_m\nabla^2J_m(\Theta^*)\bigr)^{-1}\cdot k_e^2(b_{\mathrm{frozen}} + b_{\mathrm{osc}})
      + O(k_e^{-2}).
    \end{align*}
    The offset is $O(1/k_e)$. Its orientation-odd part is the quantity measured in Remark H.7a, so the prediction
    $\Theta_+ - \Theta_- = -2J^{-1}(b_{\mathrm{odd}} + b_{\mathrm{osc}}) + O(k_e^{-2})$ of that remark is now a
    statement with error control.
\end{enumerate}
\end{thm}
\status{Proved (tube and slaving, periodic Riccati block-triangularisation, second-order averaging, Poincaré map);
checked by Newton shooting for the exact orbit, its monodromy matrix and its period mean against the averaged Hessian
and the first-order prediction, at $k_e=10$--$160$ (\code{tests/test_periodic.py}, \code{exp_periodic_orbit.py}).
Novelty (literature audit): the first step of \S8 branch 1; the tools (two-time-scale averaging under attraction) are
classical.}

\begin{proof}
Write $\varepsilon := 1/k_e$. All constants below are uniform in $\varepsilon\le1/k_0$ and in $\Theta$ in the ball
$\mathcal K := \{\norm{\Theta - \Theta^*} \le 2\rho_0\}$, because the saturated features and their derivatives are
bounded uniformly in $\Theta$ and $f\in C^2$ on the compact tube $\mathcal T_R$.

\smallskip
\emph{Step 1 (the tube and the slaved error).} Let
$M_r := \sup\{\norm{r_\Theta(x)} : x\in\mathcal T_R,\ \Theta\in\mathcal K\}$ and $L_r := \sup\norm{\partial_xr_\Theta}$
on the same set. For $k_e > M_r/R$ the tube $\{\norm{e}\le R\}$ is forward invariant: on $\norm{e} = R$,
$\frac{\mathrm{d}}{\mathrm{d}t}\frac12\norm{e}^2 = e^\T r_\Theta(x) - k_e\norm{e}^2 < 0$. Put
$w := e - \varepsilon r_\Theta(x)$ with $x$ the actual state. Since $r_\Theta(x) - k_e e = -k_e w$ and
$\dot x = \dot x_m - k_e w$,
\[
  \dot w = -\bigl(k_eI - \partial_xr_\Theta(x)\bigr)w - \varepsilon\,\partial_xr_\Theta(x)\dot x_m
  - \varepsilon\,\partial_\Theta r_\Theta(x)\cdot\dot\Theta,
\]
and $\dot\Theta = F(e,\Theta,t) := \bigl(G\nabla_\vartheta[e^\T\tilde\Phi(x;\vartheta_i)]
- \varepsilon\sigma_0(\vartheta_i - \vartheta^0)\bigr)_i$ is $O(\norm{e} + \varepsilon) = O(\norm{w} + \varepsilon)$.
So for $k_e\ge2L_r + C$, $\frac{\mathrm{d}}{\mathrm{d}t}\norm{w} \le -(k_e/2)\norm{w} + C_1\varepsilon$ on
$\{\norm{w}\le1\}\cap\{\norm{e}\le R\}$, whence $\norm{w(t)} \le e^{-k_et/2}\norm{w(0)} + 2C_1\varepsilon^2$. The set
$W_\varepsilon := \{\norm{w} \le 3C_1\varepsilon^2\}$ is forward invariant and absorbs every solution of the tube
within $t_1 = O(\varepsilon\log(1/\varepsilon))$, during which $\Theta$ moves by
$O(\varepsilon\norm{e(0)} + \varepsilon^2\log(1/\varepsilon))$ ($\dot\Theta = O(\norm{e} + \varepsilon)$ and $\norm{e}$
decays at rate $k_e/2$). On $W_\varepsilon$,
$e = \varepsilon r_\Theta(x_m(t)) + \varepsilon[r_\Theta(x_m + e) - r_\Theta(x_m)] + w
= \varepsilon r_\Theta(x_m(t)) + O(\varepsilon^2)$.

\smallskip
\emph{Step 2 (the slow equation and the fast variable).} Let $e_\Theta(t)$ be the frozen-$\Theta$ periodic error of
Proposition H.7(a): the unique $T_0$-periodic solution of $\dot e = r_\Theta(x_m(t) + e) - k_e e$, exponentially
stable with rate at least $k_e/2$, with
$e_\Theta = \varepsilon r_\Theta(x_m) + \varepsilon^2(\partial_xr_\Theta\,r_\Theta - \dot r_\Theta)(x_m) + O(\varepsilon^3)$
in sup norm (Taylor expansion of $r_\Theta$ in $x$ to second order; $f\in C^2$). Because $\Theta$ enters
$r_\Theta = f - \fhat_\Theta$ only through the smooth $\fhat_\Theta$, $e_\Theta(t)$ is $C^2$ in $\Theta$ as soon as
$f\in C^1$, and differentiating the frozen equation (gain $1/k_e$) gives
$\partial_\Theta e_\Theta = \varepsilon\,\partial_\Theta r_\Theta(x_m) + O(\varepsilon^2)$ and
$\partial^2_\Theta e_\Theta = O(\varepsilon)$, uniformly on $\mathcal K$. Let $\zeta := e - e_{\Theta(t)}(t)$ be the
fast variable. Then
\[
  \dot\zeta = \bigl[r_\Theta(x_m + e) - r_\Theta(x_m + e_\Theta)\bigr] - k_e\zeta - \partial_\Theta e_\Theta(t)\cdot\dot\Theta,
\]
where the bracket is Lipschitz in $\zeta$ with constant $L_r$ and the last term is
$O(\varepsilon)\cdot O(\varepsilon) = O(\varepsilon^2)$ ($\dot\Theta = F = O(\norm{e} + \varepsilon)$). So for
$k_e\ge2L_r + C$ the set $Z_\varepsilon := \{\norm{\zeta} \le C_2\varepsilon^3\}$ is forward invariant and absorbs
every solution of the tube within $O(\varepsilon\log(1/\varepsilon))$; on $W_\varepsilon$ of Step 1 one has
$\norm{\zeta} = O(\varepsilon^2)$ at once and $O(\varepsilon^3)$ after this further transient. On $Z_\varepsilon$,
since $F$ is $C^1$ in $e$ with $\partial_eF = O(1)$,
\begin{equation}
  \dot\Theta = F(e_\Theta(t),\Theta,t) + R(\Theta,\zeta,t)
  = \varepsilon H(\Theta,t) + \varepsilon^2H_2(\Theta,t) + O(\varepsilon^3),\qquad
  \norm{R} \le C\norm{\zeta} = O(\varepsilon^3),
  \tag{S}
\end{equation}
where $\varepsilon H + \varepsilon^2H_2$ is the expansion of the frozen-$\Theta$ drift along the frozen periodic
error, and by Proposition H.7(b),
$\langle H_2\rangle(\Theta)_i = G\nabla_\vartheta(U_{\mathrm{even}} + U_{\mathrm{odd}})[\mu_\Theta](\vartheta_i)$, that
is, $\varepsilon^2\langle H_2\rangle(\Theta^*) = b_{\mathrm{frozen}}$. $H$ and $H_2$ are $T_0$-periodic and $C^2$ in
$\Theta$ on $\mathcal K$, with bounds uniform in $\varepsilon$.

\smallskip
\emph{Regularity of the remainders (used in Steps 4 and 6).} The second-order terms of \textup{(S)}, and of
\textup{(N)} below, are used in sup norm only, for the location (d). The implicit-function and contraction arguments
of Steps 4 and 6 need first-order $C^1$ closeness, which holds: $F(e_\Theta(t),\Theta,t)$ is $C^2$ in $\Theta$ with
$\partial_\Theta[F(e_\Theta(t),\Theta,t)] = \varepsilon\,\partial_\Theta H(\Theta,t) + O(\varepsilon^2)$ (from
$\partial_\Theta e_\Theta = \varepsilon\,\partial_\Theta r_\Theta(x_m) + O(\varepsilon^2)$ and the smooth dependence of
$F$ on $(e,\Theta)$ through $\tilde\Phi$); the remainder $R$ is Lipschitz in $\zeta$ with constant
$\sup\norm{\partial_eF} = O(1)$ and Lipschitz in $\Theta$ with constant $O(\norm{\zeta}) = O(\varepsilon^3)$; and the
forcing $-\partial_\Theta e_\Theta\cdot\dot\Theta$ of the $\zeta$-equation is Lipschitz in $\Theta$ with constant
$O(\varepsilon^2)$ (from $\partial^2_\Theta e_\Theta = O(\varepsilon)$ and $\partial_\Theta F = O(\varepsilon)$). None
of this uses more than $f\in C^2$.

\smallskip
\emph{Step 3 (second-order averaging normal form).} Let $u$ be the zero-mean antiderivative of $H - \langle H\rangle$
and $\Theta = Y + \varepsilon u(Y,t)$. Then
$\dot Y = \varepsilon\langle H\rangle(Y) + \varepsilon^2\bigl[H_2(Y,t) + D_\Theta H(Y,t)u(Y,t)
- D_Yu(Y,t)\langle H\rangle(Y)\bigr] + O(\varepsilon^3)$. Let $u_2$ be the zero-mean antiderivative of the
fluctuating part of the bracket and $Y = Y_2 + \varepsilon^2u_2(Y_2,t)$. Since
$\langle D_Yu\rangle = D_Y\langle u\rangle = 0$,
\begin{equation}
  \dot Y_2 = \varepsilon\langle H\rangle(Y_2) + \varepsilon^2\Phi_2(Y_2) + O(\varepsilon^3),\qquad
  \Phi_2 := \langle H_2\rangle + \langle D_\Theta H\cdot u\rangle,
  \tag{N}
\end{equation}
the second-order averaged equation of Remark H.7a (Sanders--Verhulst--Murdock~\cite[\S2.9]{svm2007}), now with a
remainder that is uniformly $O(\varepsilon^3)$ on $Z_\varepsilon\times\mathcal K$ because \textup{(S)} holds there.
The transformation $\Theta\mapsto Y_2$ is a $T_0$-periodic near-identity diffeomorphism from $\mathcal K$ onto its
image, $O(\varepsilon)$-close to the identity with $O(\varepsilon)$ derivative, and the period mean of $\Theta$ along
any curve $Y_2(t) = Y_2^0 + O(\varepsilon^2)$ equals the period mean of $Y_2$ up to $O(\varepsilon^3)$, because $u$
and $u_2$ have zero mean.

\smallskip
\emph{Step 4 (the Poincaré map in the slow coordinates).} Let $P = (P_Y,P_\zeta)$ be the time-$T_0$ map of
\textup{(CL)} from $t = 0$, written in the coordinates $(Y_2,\zeta)$ on $\mathcal K\times Z_\varepsilon$. By
\textup{(N)} and Step 2,
\begin{gather*}
  P_Y(Y_2,\zeta) = Y_2 + \varepsilon T_0\tilde A(Y_2) + O(\varepsilon^3),\qquad
  \tilde A := \langle H\rangle + \varepsilon\bigl[\Phi_2 + (T_0/2)D\langle H\rangle\cdot\langle H\rangle\bigr],\\
  \norm{P_\zeta(Y_2,\zeta)} \le e^{-k_eT_0/2}\norm{\zeta} + O(\varepsilon^3),
\end{gather*}
the $\varepsilon^2(T_0^2/2)D\langle H\rangle\langle H\rangle$ term being the second-order term of the time-$T_0$ flow
of \textup{(N)}. The Lipschitz constants of the coupling are $\Lip_\zeta P_Y = O(\varepsilon)$ (a perturbation of
$\zeta$ acts on $\dot\Theta$ through $\partial_eF = O(1)$ while it decays at rate $k_e/2$, so its integrated effect
over a period is $O(\varepsilon)\norm{\zeta}$) and $\Lip_YP_\zeta = O(\varepsilon^3)$ (the forcing of the
$\zeta$-equation is $O(\varepsilon^2)$-Lipschitz in $\Theta$ and is integrated against the fast decay). Hence
$\zeta\mapsto P_\zeta(Y_2,\zeta)$ has a unique fixed point $\zeta^*(Y_2) = O(\varepsilon^3)$, Lipschitz in $Y_2$ with
constant $O(\varepsilon^3)$, and the fixed points of $P$ are the zeros of
\[
  \tilde A_\varepsilon(Y) := \frac{P_Y(Y,\zeta^*(Y)) - Y}{\varepsilon T_0} = \tilde A(Y) + O(\varepsilon^2).
\]
By the regularity paragraph, $D_YP_Y = I + \varepsilon T_0D\langle H\rangle(Y) + O(\varepsilon^2)$ (the variational
equation of the slow flow over one period, whose coefficient is $\varepsilon\,\partial_\Theta H + O(\varepsilon^2)$),
and the $\zeta^*$-dependence contributes $\Lip_\zeta P_Y\cdot\Lip_Y\zeta^* = O(\varepsilon^4)$, so
$\tilde A_\varepsilon$ is Lipschitz with $\Lip(\tilde A_\varepsilon - \langle H\rangle) = O(\varepsilon)$ on
$\mathcal K$. Since $D\langle H\rangle(\Theta^*) = -G_m\nabla^2J_m(\Theta^*)$ is invertible,
$Y\mapsto Y - D\langle H\rangle(\Theta^*)^{-1}\tilde A_\varepsilon(Y)$ maps a ball of $\varepsilon$-independent radius
$\rho_0$ around $\Theta^*$ into itself and is a contraction there for $\varepsilon$ small (the Lipschitz form of the
implicit function theorem), which gives a unique zero $Y_2^* = \Theta^* + O(\varepsilon)$, hence a unique fixed point
of $P$ with $\Theta$-component within $O(\varepsilon)$ of $\Theta^*$, that is, a unique $T_0$-periodic solution in the
stated set; the bound on $x_p - x_m$ is Step 1. Expanding $\tilde A_\varepsilon(Y_2^*) = 0$ to first order,
$D\langle H\rangle(\Theta^*)(Y_2^* - \Theta^*) + \varepsilon\Phi_2(\Theta^*) = O(\varepsilon^2)$ (the term
$(T_0/2)D\langle H\rangle\langle H\rangle(Y_2^*)$ is $O(\varepsilon)$ because $\langle H\rangle(Y_2^*) = O(\varepsilon)$),
so
\[
  Y_2^* = \Theta^* - \varepsilon D\langle H\rangle(\Theta^*)^{-1}\Phi_2(\Theta^*) + O(\varepsilon^2)
  = \Theta^* - J^{-1}(b_{\mathrm{frozen}} + b_{\mathrm{osc}}) + O(\varepsilon^2),
\]
using $J = \varepsilon D\langle H\rangle(\Theta^*)$ and $\varepsilon^2\Phi_2(\Theta^*) = b_{\mathrm{frozen}} + b_{\mathrm{osc}}$.
Along the periodic orbit $Y_2(t) = Y_2^* + O(\varepsilon^2)$ (since $\langle H\rangle(Y_2^*) = O(\varepsilon)$), so by
Step 3 the period mean of $\Theta_p$ is $Y_2^* + O(\varepsilon^2)$: this is \textup{(d)}, and
$\Theta_p(t) - \bar\Theta_p = \varepsilon u(\Theta^*,t) + O(\varepsilon^2)$ is \textup{(a)}.

\smallskip
\emph{Step 5 (the monodromy matrix).} Linearise \textup{(CL)} along the periodic orbit: with $\delta e$ and
$\delta\Theta$,
\[
  \delta\dot e = (\partial_xr - k_eI)\delta e + \partial_\Theta r\cdot\delta\Theta,\qquad
  \delta\dot\Theta = \partial_eF\cdot\delta e + \partial_\Theta F\cdot\delta\Theta,
\]
all coefficients $T_0$-periodic and evaluated along the orbit; $\partial_xr$, $\partial_\Theta r$, $\partial_eF$ are
$O(1)$ and $\partial_\Theta F = O(\varepsilon)$ ($F$ is linear in $e$ with coefficients smooth in $\Theta$, plus the
leak $-\varepsilon\sigma_0(\vartheta_i - \vartheta^0)$, so $\partial_\Theta F = O(\norm{e} + \varepsilon)$, and
$e = O(\varepsilon)$ on the orbit). Let $S(t)$ be a $T_0$-periodic solution of the Riccati equation
\[
  \dot S = (\partial_xr - k_eI)S + \partial_\Theta r - S\,\partial_eF\,S - S\,\partial_\Theta F.
\]
For $k_e$ large it exists, is unique with $\norm{S} \le C/k_e$, and depends smoothly on the data (contraction of the
periodic solution operator of the dominant part $\dot S = -k_eS + g$; cf.\ Proposition H.7(a)). In the coordinates
$\eta := \delta e - S\delta\Theta$ the linearisation is block triangular:
\[
  \dot\eta = \bigl[(\partial_xr - k_eI) - S\,\partial_eF\bigr]\eta,\qquad
  \delta\dot\Theta = \partial_eF\cdot\eta + \bigl[\partial_eF\,S + \partial_\Theta F\bigr]\delta\Theta,
\]
so the monodromy matrix is similar to a block-triangular matrix whose diagonal blocks are the monodromies of the two
diagonal equations. The $\eta$-block has norm at most $e^{-(k_e - L)T_0}$ with
$L := L_r + \sup\norm{S\,\partial_eF} = O(1)$ (Grönwall): these are the $n$ error-channel multipliers. The
$\delta\Theta$-block is the linearisation along the orbit of the slow equation \textup{(S)} in $\Theta$-coordinates,
$\varepsilon D_\Theta H(\Theta_p(t),t) + O(\varepsilon^2) = \varepsilon D_\Theta H(\Theta^*,t) + O(\varepsilon^2)$, whose
monodromy is $I + \varepsilon T_0D\langle H\rangle(\Theta^*) + O(\varepsilon^2)
= I - \varepsilon T_0G_m\nabla^2J_m(\Theta^*) + O(\varepsilon^2)$. Its eigenvalues are
$1 - \varepsilon T_0\nu_j + O(\varepsilon^2) = \exp(-\varepsilon T_0\nu_j + O(\varepsilon^2))$, those of
$G_m\nabla^2J_m(\Theta^*)$ being the $\nu_j$ by similarity with $G_m^{1/2}\nabla^2J_m(\Theta^*)G_m^{1/2}$; the
$O(\varepsilon^2)$ perturbation moves them by $O(\varepsilon^2)$ because $I - \varepsilon T_0G_m\nabla^2J_m(\Theta^*)$
is diagonalisable (similar to a symmetric matrix through $G_m^{1/2}$), so the Bauer--Fike theorem applies with a
condition number independent of $\varepsilon$ --- this covers the repeated eigenvalues of a collapsed configuration.
This is \textup{(b)}.

\smallskip
\emph{Step 6 (the basin).} Let $\nabla^2J_m(\Theta^*)\succ0$. In the norm $\norm{v}_{\Ginv} := \norm{G_m^{-1/2}v}$,
the matrix $I - \varepsilon T_0G_m\nabla^2J_m(\Theta^*)$ is similar to a symmetric contraction with norm exactly
$1 - \varepsilon T_0\lambda_*$ once $\varepsilon T_0\nu_{mp} < 1$. By Step 4,
$D_YP_Y = I - \varepsilon T_0G_m\nabla^2J_m(Y) + O(\varepsilon^2)$ on $\mathcal K$, and $\nabla^2J_m$ is Lipschitz, so
on the ball $\{\norm{Y_2 - Y_2^*}_{\Ginv} \le \rho_0\}$ the Lipschitz constant of $P_Y$ in $Y_2$ is at most
$1 - \varepsilon T_0\lambda_*(1 - \kappa)$ with $\kappa = O(\rho_0 + \varepsilon)$; the coupling constants are
$\Lip_\zeta P_Y = O(\varepsilon)$ and $\Lip_YP_\zeta = O(\varepsilon^3)$, and $P_\zeta$ contracts $\zeta$ by
$e^{-k_eT_0/2}$ (Step 4). The $2\times2$ matrix of Lipschitz constants
\[
  \begin{pmatrix} 1 - \varepsilon T_0\lambda_*(1 - \kappa) & C\varepsilon\\ C\varepsilon^3 & e^{-k_eT_0/2}\end{pmatrix}
\]
has spectral radius $1 - \varepsilon T_0\lambda_*(1 - \kappa')$ with $\kappa' = O(\rho_0 + \varepsilon)$, so on
$\{\norm{Y_2 - Y_2^*}_{\Ginv} \le \rho_0\}\times Z_\varepsilon$, with $\rho_0$ small but independent of
$\varepsilon$, $P$ is a contraction in an adapted product norm with that factor: its iterates converge to the fixed
point geometrically, uniformly in $\varepsilon$. Solutions started in the tube with $\norm{e(0) - e_p(0)} \le 1$ enter
$Z_\varepsilon$ within $O(\varepsilon\log(1/\varepsilon))$, during which $\Theta$ moves by
$O(\varepsilon)\norm{e(0) - e_p(0)}$ (Step 1), which is absorbed by halving $\rho_0$. Between the sampling times the
solutions stay within a factor $C$ of each other (Grönwall on one period, with the fast variable absorbing
$\norm{e(0) - e_p(0)}$ into an $O(\varepsilon)\norm{e(0) - e_p(0)}$ displacement of $\Theta$, Step 1). Transporting
back through the near-identity transformations gives \textup{(c)}.
\end{proof}

\paragraph{Numerics} (\code{exp_periodic_orbit.py} $\to$ \code{results/periodic_orbit/table.md};
\code{tests/test_periodic.py}). Two instances, $k_e = 10$, $20$, $40$, $80$, $160$. For each $k_e$ the periodic orbit
of (CL) is found by Newton shooting on the Poincaré map (DOP853 at relative tolerance $10^{-11}$; residual below
$10^{-11}$ in every case), its monodromy matrix is the finite-difference Jacobian of the map, and its period mean is
compared with $\Theta^*$ and with the first-order prediction $\Theta^* - J^{-1}(b_{\mathrm{frozen}} + b_{\mathrm{osc}})$
of (d). Distances and cosines are in the labelled $G^{-1}$ metric.
\begin{itemize}
  \item \textbf{Instance A} ($n = 2$, $m = 3$, the plant and elliptic reference of Remark H.7a, $\sigma_0 = 0.5$,
    $G = \operatorname{diag}(4,4,1,1,1)$). The minimiser of $J_m$ is a collapsed configuration (kernel spread
    $4\cdot10^{-14}$) with $\lambda_{\min}(\nabla^2J_m) = 0.0148$: a non-degenerate minimum with permutation symmetry.
    \begin{center}
    \small
    \begin{tabularx}{\linewidth}{@{}llYYY@{}}
    \toprule
    $k_e$ & $\norm{\bar\Theta_p - \Theta^*}$ & full prediction: size ratio, cosine, vector error
      & frozen-measure part only: vector error
      & parameter Floquet exponents against $(1/k_e)\nu_j$: max, median relative error\\
    \midrule
    10  & $4.09\cdot10^{-2}$ & 1.149, 0.9972, 0.169 & 0.398 & $8.3\cdot10^{-2}$, $1.2\cdot10^{-2}$\\
    20  & $1.91\cdot10^{-2}$ & 1.074, 0.9993, 0.083 & 0.360 & $7.3\cdot10^{-2}$, $7.3\cdot10^{-3}$\\
    40  & $9.22\cdot10^{-3}$ & 1.036, 0.9998, 0.041 & 0.348 & $4.5\cdot10^{-2}$, $4.0\cdot10^{-3}$\\
    80  & $4.53\cdot10^{-3}$ & 1.018, 1.0000, 0.020 & 0.343 & $2.5\cdot10^{-2}$, $2.1\cdot10^{-3}$\\
    160 & $2.24\cdot10^{-3}$ & 1.009, 1.0000, 0.010 & 0.342 & $1.3\cdot10^{-2}$, $1.1\cdot10^{-3}$\\
    \bottomrule
    \end{tabularx}
    \end{center}
    The offset scales like $k_e^{-1.05}$ (fitted exponent; (a) predicts $-1$), the vector error of the full
    prediction halves with every doubling of $k_e$ (the $O(k_e^{-2})$ remainder of (d) against an $O(1/k_e)$ offset),
    and the frozen-measure part alone has a 35\% vector error at every $k_e$: the oscillation term $b_{\mathrm{osc}}$
    of Remark H.7a is an $O(1)$ fraction of the first-order location. The median relative error of the Floquet
    exponents is $O(1/k_e)$, as (b) states; the largest relative error sits on the $(m-1)$-fold degenerate separation
    pair with the smallest exponent ($0.0018$ at $k_e = 20$), where an $O(k_e^{-2})$ absolute error is a large
    fraction of a small exponent. The two error-channel multipliers are below the resolution of the finite-difference
    monodromy (exponents of order $k_e$); their product is recovered from Liouville's identity
    $\sum\log\abs{\lambda} = \int_0^{T_0}\tr\,\partial F/\partial z\dd t$ along the orbit, which gives a mean
    error-channel exponent of $k_e - 0.16$, $k_e - 0.11$, $k_e - 0.086$, $k_e - 0.074$, $k_e - 0.068$ at
    $k_e = 10\dots160$ (instance B: $k_e + 0.05\dots k_e + 0.09$): the multipliers are $e^{-(k_e - L)T_0}$ with
    $\abs{L} < 0.2$, as (b) states.
  \item \textbf{Instance B} ($n = 1$, $m = 2$, a two-bump plant
    $f(x) = 1.6e^{-(x-1)^2/0.4} - 1.2e^{-(x+1)^2/0.4} + 0.1x$ on $x_m = 2\sin t$, $\sigma_0 = 0.05$,
    $G = \operatorname{diag}(2,1,1)$). The minimiser has two distinct kernels of opposite sign (spread $2.1$,
    $\lambda_{\min}(\nabla^2J_m) = 0.066$). Here $b_{\mathrm{osc}} = 0$ exactly: for a one-dimensional sinusoidal
    sweep the reflection $t\mapsto\pi - t$ maps $x_m$ to itself and $u$ to $-u$, so $\langle D_\Theta H\cdot u\rangle = 0$,
    and the full and frozen predictions coincide. The reference is fast on the kernel scale ($\sup\norm{\dot x_m}/s$
    is $2/s$: between $1.3$ and $10$ over the admissible widths, about $2.6$ at the minimiser's width; the residual
    bandwidth is of the same order), so $k_e = 10$--$20$ is outside the averaging regime:
    \begin{center}
    \small
    \begin{tabularx}{\linewidth}{@{}llYY@{}}
    \toprule
    $k_e$ & $\norm{\bar\Theta_p - \Theta^*}$ & prediction: size ratio, cosine, vector error
      & Floquet: max, median relative error\\
    \midrule
    10  & $6.44\cdot10^{-2}$ & 2.818, 0.880, 1.995 & $9.5\cdot10^{-2}$, $4.0\cdot10^{-2}$\\
    20  & $2.22\cdot10^{-2}$ & 1.942, 0.929, 1.078 & $3.2\cdot10^{-2}$, $1.4\cdot10^{-2}$\\
    40  & $8.37\cdot10^{-3}$ & 1.464, 0.970, 0.551 & $1.2\cdot10^{-2}$, $5.4\cdot10^{-3}$\\
    80  & $3.51\cdot10^{-3}$ & 1.227, 0.990, 0.277 & $4.6\cdot10^{-3}$, $2.2\cdot10^{-3}$\\
    160 & $1.59\cdot10^{-3}$ & 1.111, 0.9969, 0.139 & $2.0\cdot10^{-3}$, $9.5\cdot10^{-4}$\\
    \bottomrule
    \end{tabularx}
    \end{center}
    The vector error and the Floquet error halve with each doubling of $k_e$, as the theorem's $O(k_e^{-2})$
    remainders predict, but the constants are large: the asymptotic regime starts around $k_e\approx80$ here. This
    instance shows what the theorem does not claim: the $O(1/k_e)$ location is sharp only once $k_e$ exceeds the
    effective bandwidth $\omega_{\mathrm{eff}}$ of Proposition G's (H2a$'$) (Section~\ref{sec:propG}).
\end{itemize}
\checkedby{\code{tests/test_periodic.py}: $\nabla J_m$ against finite differences ($10^{-7}$) and the non-degenerate
two-kernel minimum of instance B; at $k_e = 20$ and $40$ the shooting residual is below $10^{-10}$,
$\norm{\bar\Theta_p - \Theta^*}$ and $\max\norm{e}$ fall by more than 35\% from $20$ to $40$, the vector error of the
first-order location falls by more than 40\% and is below $0.6$ at $k_e = 40$, and the maximal relative Floquet error
falls by more than 40\% and is below 2\% at $k_e = 40$.}

\begin{rmk}{Remark}{J.1}[what Theorem J says about the program of Section~\ref{sec:program}]
\leavevmode
\begin{itemize}
  \item It is the deterministic, finite-$m$ first step of branch 1: near every hyperbolic critical point of the
    $m$-particle regularised loss $k_e\cdot\mathcal G$ at $\tau = 0$ the real closed loop has exactly one stationary
    state (a periodic orbit), an $O(1/k_e)$-perturbation of it with the same stability type, and the convergence rate
    is the one of the averaged gradient flow in slow time. (The theorem does not say that every periodic orbit of
    (CL) arises this way.) The two missing pieces of Conjecture H.6 remain: dither ($\tau>0$, where the stationary
    state is a measure) and uniformity in $m$.
  \item It makes the stationary (period-mean) prediction of Remark H.7a a theorem for this system: the first-order
    location of the orbit, hence the orientation difference of the period-mean kernels, is
    $\Theta^* - J^{-1}(b_{\mathrm{frozen}} + b_{\mathrm{osc}})$ with an $O(k_e^{-2})$ error. When the leading offset
    is itself small, the $O(k_e^{-2})$ remainder can dominate its \emph{direction} at moderate $k_e$: a collapsed
    saddle instance examined in the round-5 review (\code{review/REVIEW_round5_J.md}) has a first-order location
    nearly orthogonal to the true offset at $k_e = 20$ and aligned with it at $k_e = 80$, while the sizes obey the
    theorem's orders throughout.
  \item The leak is not needed for the structure, only for the hyperbolicity of $\Theta^*$. At a collapsed critical
    point the separation Hessian is $\sigma_0G^{-1} - R_2$ (Proposition K, Section~\ref{sec:thmK}); for $n = 1$ and
    $\sigma_0 = 0$ the amplitude row of $R_2$ vanishes at a leak-free critical point
    ($\langle r\varphi\rangle = \partial_c\langle r\varphi\rangle = \partial_s\langle r\varphi\rangle = 0$ there), so
    the $m-1$ amplitude-redistribution directions are exactly flat and the theorem does not apply, while for $n\ge2$
    only the amplitude block of $R_2$ vanishes and a collapsed critical point is generically non-degenerate even
    without a leak. Where $\nabla^2J_m$ is singular the theorem does not apply.
  \item At a collapsed critical point (all $\vartheta^*_i$ equal) the permutation symmetry makes the separation modes
    $(m-1)$-fold degenerate; the Floquet exponents come in such multiplets (\code{results/periodic_orbit/table.md});
    Proposition K (Section~\ref{sec:thmK}) works out this case.
\end{itemize}
\end{rmk}
